\section{Related Work}

\paragraph{Vision-language models.}
Contrastive image-text pretraining and multimodal instruction tuning have produced strong general-purpose visual representations and assistants \citep{radford2021clip,jia2021align,zhai2022lit,alayrac2022flamingo,wang2022ofa,chen2023pali,wang2023beit3,zhai2023siglip,li2022blip,li2023blip2,dai2023instructblip,liu2023llava,lu2023unifiedio}. These models provide the semantic layer needed for open-vocabulary monitoring, but they are usually evaluated with preselected visual inputs or bounded single-query contexts. Our work treats VLM access itself as a scarce action and asks when a monitoring agent should spend semantic compute.

\paragraph{Video and long-video understanding.}
Spatiotemporal CNNs, video transformers, and masked video learners improved action recognition and video representation learning \citep{tran2015c3d,carreira2017quo,feichtenhofer2019slowfast,bertasius2021timesformer,arnab2021vivit,liu2022videoswin,tong2022videomae,tong2023videomaev2}. Video-language pretraining and temporal grounding methods connect these representations to language \citep{lei2021clipbert,bain2021frozen,huang2024vtimellm}. Long-video VLMs and memory-based systems address context length, temporal grounding, and streaming interaction \citep{song2024moviechat,chen2024videollmonline,mangalam2023egoschema,fu2025videomme}. These works typically focus on one video or a small set of videos. Multi-stream monitoring additionally requires cross-stream allocation under a real-time budget.

\paragraph{Video benchmarks and event-centric evaluation.}
Large-scale video benchmarks have shaped action recognition, activity localization, dense captioning, egocentric memory, and spatiotemporal action detection \citep{caba2015activitynet,krishna2017dense,sigurdsson2016hollywood,goyal2017something,gu2018ava,grauman2022ego4d}. They define important semantic and temporal capabilities, but they do not directly measure how an agent should distribute expensive semantic queries across hundreds of simultaneous streams. \benchmark builds on their event-centric view while adding budgeted VLM calls, false alarms per hour, time-to-detect, and cross-stream contention.

\paragraph{Efficient perception.}
Dynamic networks, early exits, adaptive frame selection, token sparsification, token merging, and adaptive token budgets reduce inference cost within a selected input \citep{figurnov2017spatially,huang2018msdnet,wang2018skipnet,wu2019adaframe,rao2021dynamicvit,meng2022evit,yin2022avit,ryoo2021tokenlearner,bolya2023tome}. Such methods are complementary to \method: they make a selected observation cheaper, while our scheduler decides which observations should be selected in the first place. This distinction matters in monitoring because the dominant cost is often not a single forward pass, but the opportunity cost of looking at one stream while ignoring many others.

\paragraph{Budgeted information acquisition.}
Adaptive submodularity and sensor placement formulate partial observability as sequential information acquisition under a budget \citep{krause2008sensor,golovin2011adaptive}. Our problem is related, but the observation action is a VLM query whose value depends on event rarity, stream state, prompt semantics, temporal delay, and downstream alert utility. This setting motivates the value-of-information view in \method rather than a purely dense-recognition or fixed-sampling protocol.

\paragraph{Video anomaly detection.}
Surveillance anomaly detection benchmarks and methods study rare events, weak supervision, multimodal violence detection, normality modeling, and abnormality scoring \citep{hasan2016regularity,delgiorno2016discriminative,luo2017revisit,sultani2018ucfcrime,liu2018futureframe,park2020memoryguided,wu2020xdviolence,georgescu2021selfsupervised,tian2021rtfm,feng2021mist}. Their datasets are a natural substrate for monitoring, but the standard metrics often omit VLM-call budgets, natural-language incident reports, and multi-stream scheduling.

\paragraph{Agentic systems and tool use.}
Language agents have been studied through embodied and web environments, tool use, reasoning-action loops, API calling, and search-based planning \citep{shridhar2021alfworld,yao2022webshop,yao2023react,schick2023toolformer,yao2023tot,deng2023mind2web,qin2024toolllm,liu2024agentbench,zhou2024webarena}. We reinterpret visual observation as an agent action: the agent can spend compute to look at a stream, ask a visual question, update memory, or defer. This exposes perception allocation as a first-class decision problem for multimodal agents operating under a hard real-time budget.
